% Generated draft: requires mathematical and source review.
\documentclass[11pt]{article}
\usepackage[margin=25mm]{geometry}
\usepackage{fontspec}
\setmainfont{Latin Modern Roman}
\usepackage{amsmath,amssymb,mathtools,unicode-math}
\setmathfont{Latin Modern Math}
\usepackage{enumitem,needspace,xurl,hyperref}
\setlength{\parindent}{0pt}
\setlength{\parskip}{5pt}
\setlength{\emergencystretch}{3em}
\newcommand{\N}{\mathbb{N}}
\newcommand{\Z}{\mathbb{Z}}
\newcommand{\Q}{\mathbb{Q}}
\newcommand{\R}{\mathbb{R}}
\newcommand{\C}{\mathbb{C}}
\newcommand{\F}{\mathbb{F}}
\newcommand{\A}{\mathbb{A}}
\newcommand{\PP}{\mathbb{P}}
\newcommand{\E}{\mathbb{E}}
\newcommand{\eps}{\varepsilon}
\newcommand{\dd}{\,\mathrm{d}}
\begin{document}
\section*{UnsolvedMath ID: 20001124}
\textbf{Source problem code:} catalogue entry
\subsection{Definitions and mathematical statement}
Let \(n\ge2\) be an integer, let \(\mathcal L\) be a set of \(n^2\) distinct affine lines in \(\R^3\), and let \(S\subset\R^3\) be finite. Interpret ``each containing \(n\) points'' as
\[
 |\ell\cap S|\ge n\qquad(\ell\in\mathcal L).
\]
Assume no \(n\) lines are coplanar, equivalently each affine plane contains at most \(n-1\) members of \(\mathcal L\). J. Bourgain's source problem asks for a lower bound on \(|S|\). Its OCR example ``\(n3-\le\)'' is interpreted by the aid as a request of the form \(n^{3-\varepsilon}\); the exact exponent typography is not independently recovered here.

The supplied aid reports that the published Guth--Katz rich-line theorem gives the stronger sharp result: there is an absolute \(c>0\) such that every configuration satisfying these hypotheses has
\[
 |S|\ge cn^3.
\]
The constant is independent of \(n,S,\mathcal L\). Sharpness here means the power \(3\), not an exact leading constant. The aid presents this as a supplied literature resolution, requiring verification of the cited theorem rather than a new proof.

A separate reported construction varies the coplanarity cap. For integers
\[
 n\ge2,\qquad n\le s\le n^2,\qquad s\mid n^2,
\]
the aid asserts a configuration of exactly \(n^2\) distinct lines and a finite point set \(S_{n,s}\) such that every line has \(n\) selected incident points, every plane contains at most \(s\) lines, and
\[
 |S_{n,s}|\le4n^{7/2}s^{-1/2}.
\]
``Selected'' does not assert that the full set has exactly \(n\) incidences on each line. The proof is said to use designated parallel block planes while also controlling every other plane, but coordinates and proof artifacts are absent. Since \(s\ge n\), this construction has a relaxed cap and is not automatically an example satisfying the original strict \(n-1\) cap, even when \(s=n\). It is retained as an additional parameterized result, not conflated with the original theorem.
\subsection{Short English statement}
Bourgain's problem asks how many points are needed for \(n^2\) lines in three-space, each rich in points, when fewer than \(n\) lines lie in any plane. The supplied aid reports the sharp cubic lower bound from Guth--Katz and separately gives a construction for larger allowed plane caps.
\subsection{Sources}
\begin{itemize}
\item Local UnsolvedMath registry, permanent ID \(20001124\), source code \texttt{catalogue entry}; \texttt{aim-combinatorics-notes.json}, \texttt{notes[248]}, tag \texttt{problem}.
\item AIM workshop ``Recent trends in additive combinatorics,'' Problem 3.12, J. Bourgain: \url{https://aimath.org/WWN/additivecomb/additivecomb.pdf}. Statement status \texttt{exact}, referring to the stored OCR extraction. The Szemer\'edi--Trotter theorem is named in the source motivation. No PDF was accessed for this draft.
\item Guth--Katz published rich-line theorem, named by the aid as giving the sharp lower bound. No full title, date, theorem number, or URL is supplied; the additional parameterized construction has no separate authoritative citation.
\item Research attempt 1: \texttt{full\_solution}, problem status at run \texttt{solved}, rigor label \texttt{proved}, novelty confidence low; registry status \texttt{solved}, classification \texttt{SOLVED-IN-LITERATURE}. Created and updated 14 August 2026. No original remarks or literature field are supplied. This machine-generated aid is not peer reviewed; independently verify before citation.
\item Ulam AI / UnsolvedMath Contributors; CC BY 4.0.
\end{itemize}
\end{document}
